\documentclass[11pt]{article}
\usepackage[T1]{fontenc}
\usepackage[utf8]{inputenc}
\usepackage{lmodern}
\usepackage[margin=1in,headheight=15pt]{geometry}
\usepackage{amsmath,amssymb,amsthm,mathtools,mathrsfs}
\usepackage{microtype,booktabs,array,longtable,enumitem}
\usepackage{xcolor}
\usepackage[colorlinks=true,linkcolor=blue!45!black,citecolor=green!35!black,urlcolor=blue!45!black]{hyperref}
\usepackage[nameinlink,noabbrev]{cleveref}
\usepackage{fancyhdr}
\usepackage{tocloft}
\setlength{\cftbeforesecskip}{3pt}
\setlength{\cftbeforesubsecskip}{0pt}
\setlength{\cftbeforetoctitleskip}{0pt}
\setlength{\cftaftertoctitleskip}{12pt}
\pagestyle{fancy}
\fancyhf{}
\fancyhead[L]{\small Finite geometry of Hahn-coherent zeros}
\fancyhead[R]{\small\thepage}
\renewcommand{\headrulewidth}{0.3pt}
\setlength{\parskip}{0.25em}
\setlength{\emergencystretch}{2em}
\setlist[enumerate]{itemsep=0.25em,topsep=0.4em}
\newtheorem{theorem}{Theorem}[section]
\newtheorem{lemma}[theorem]{Lemma}
\newtheorem{proposition}[theorem]{Proposition}
\newtheorem{corollary}[theorem]{Corollary}
\theoremstyle{definition}
\newtheorem{definition}[theorem]{Definition}
\newtheorem{example}[theorem]{Example}
\newtheorem{remark}[theorem]{Remark}
\newtheorem{question}[theorem]{Question}
\newcommand{\C}{\mathbb C}
\newcommand{\R}{\mathbb R}
\newcommand{\Q}{\mathbb Q}
\newcommand{\N}{\mathbb N}
\newcommand{\No}{\mathbf{No}}
\newcommand{\SC}{\mathbf{No}[i]}
\newcommand{\OO}{\mathcal O}
\newcommand{\HH}{\mathcal H}
\newcommand{\RR}{\mathscr R}
\newcommand{\mm}{\mathfrak m}
\newcommand{\pp}{\mathfrak p}
\newcommand{\supp}{\operatorname{supp}}
\newcommand{\st}{\operatorname{st}}
\newcommand{\Spec}{\operatorname{Spec}}
\newcommand{\Max}{\operatorname{Max}}
\newcommand{\Tr}{\operatorname{Tr}}
\newcommand{\Res}{\operatorname{Res}}
\newcommand{\id}{\operatorname{id}}
\newcommand{\im}{\operatorname{im}}
\newcommand{\rank}{\operatorname{rank}}
\newcommand{\length}{\operatorname{length}}
\newcommand{\Frac}{\operatorname{Frac}}
\newcommand{\disc}{\operatorname{disc}}
\newcommand{\HInt}{\mathop{\int}\nolimits^{\!H}}
\newcommand{\sh}{^{\sharp}}
\newcommand{\dz}{\,dz_1\wedge\cdots\wedge dz_n}
\newcommand{\K}{K_\Gamma}
\newcommand{\val}{v}
\newcommand{\Tmon}{T^{*}}
\newcommand{\bvec}{\boldsymbol b}
\newcommand{\eps}{\varepsilon}
\newcolumntype{L}[1]{>{\raggedright\arraybackslash}p{#1}}
\hypersetup{pdftitle={Finite Geometry of Hahn-Coherent Surcomplex Zeros},pdfauthor={},pdfsubject={Support-controlled multiplicity, normal forms, finite mappings, and residues}}
\title{\textbf{Finite Geometry of Hahn-Coherent\\ Surcomplex Zeros}\\[0.6em]
\large Conservation of Multiplicity, Normal Forms,\\ and Collision-Stable Residues}
\author{Research article developed from the supplied surcomplex-analysis manuscript}
\date{September 21, 2026}
\begin{document}
\maketitle
\begin{abstract}
We develop the finite-dimensional geometry of several-variable Hahn-coherent functions near an ordinary point. The main result is a support-explicit conservation theorem: if an ordinary holomorphic complete intersection has an isolated zero of multiplicity $\mu$, then every perturbation whose Hahn support is strictly positive has a quotient algebra of dimension exactly $\mu$, with the same prescribed polynomial basis. Normal forms, division certificates, and multiplication matrices are constructed on one fixed ordinary neighborhood; their supports lie in the original support plus the positive monoid generated by the perturbation. The proof uses an explicit conjugacy of Koszul differentials, not convergence of a Newton sequence. It applies to arbitrary-rank value groups and to supports with bounded accumulation cuts.

We also prove preparation with holomorphic parameters, a monad Nullstellensatz, Noether normalization and finite-map statements for the common-domain germ ring. A multivariate Hahn contour functional descends to a perfect residue pairing on the deformed algebra, satisfies a trace--Jacobian identity, and equals a sum over the actual surcomplex roots when these are simple. This pairing remains nondegenerate when roots collide. A two-equation example of multiplicity five gives explicit multiplication matrices, the collision equation, three scale regimes, and exact residue identities.

Multivariable Hahn--Weierstrass theory and analytic Nullstellensatz results have substantial published predecessors; they are not claimed here as new discoveries. The proposed contribution is the combined, common-domain, arbitrary-support normal-form and residue package, with its explicit support bounds and workspace-independence theorem. The exact formulation was not located in the literature examined, but no exhaustive priority claim is made.
\end{abstract}
\vspace{0.5em}
\noindent\textbf{Keywords.} Surcomplex analysis; Hahn series; isolated complete intersections; Koszul complex; conservation of number; Grothendieck residue; finite analytic maps.
\clearpage
\begingroup
\fontsize{9.6}{11.3}\selectfont
\setlength{\parskip}{0pt}
\tableofcontents
\endgroup
\clearpage

\section{The problem and the contribution}\label{sec:intro}
\subsection{From one equation to an infinitesimal analytic space}
The supplied manuscript \cite{Source} constructs a coherent analytic category on thickenings of ordinary complex domains. Its coefficients are ordinary holomorphic functions, but their values are assembled into surcomplex numbers by strongly summable Hahn expansions. In one variable, its preparation theorem turns an ordinary zero of order $m$ into an exact cluster of $m$ surcomplex zeros, counted with multiplicity. Its concluding research directions ask for several-variable division and finite-mapping theorems with support control.

The several-variable problem is not merely to solve a system with nonsingular Jacobian. At a singular isolated zero, an arbitrarily small perturbation can produce several points, different infinitesimal separation scales, and nonreduced components. A correct extension must retain the entire finite algebra, rather than a selected root or a formal first-order approximation. It must also address a characteristic difficulty of the supplied framework: infinitely many ordinary coefficient functions must remain holomorphic on \emph{one common neighborhood}. Repeatedly shrinking a neighborhood at each Hahn exponent does not establish that requirement.

We answer the following precise version of this problem. Let
\[
 f=(f_1,\ldots,f_n),\qquad
 \mu=\dim_\C\C\{z_1,\ldots,z_n\}/(f_1,\ldots,f_n)<\infty,
\]
where every $f_i(0)=0$. Replace $f_i$ by
\begin{equation}\label{eq:mainperturbation}
 F_i=f_i+E_i,\qquad
 E_i=\sum_{\gamma>0}e_{i,\gamma}(z)t^\gamma,
\end{equation}
with well-ordered supports and holomorphic coefficients on a common ordinary neighborhood. Does the entire infinitesimal zero cluster still have length $\mu$? Can one construct its multiplication matrices, identify every root, and attach a residue functional that survives collisions?

The answer is affirmative in the common-domain Hahn category. The quantitative algebraic conclusion is especially simple. If $T$ is the union of the positive perturbation supports and $S$ is the support of an input $G$, all normal-form and division outputs have support contained in
\begin{equation}\label{eq:masterbound}
 S+T^*,\qquad
 T^*=\{\gamma_1+\cdots+\gamma_k:k\geq0,\ \gamma_j\in T\}.
\end{equation}
This is a support bound, not a real norm estimate and not a claim that successive approximations converge in the full surreal fine topology.

\subsection{Main statements in a fixed workspace}
Fix a nonzero divisible ordered abelian subgroup $\Gamma\subset\No$ and put
\[
 K=K_\Gamma=\C((t^\Gamma)),\qquad
 \mm_K=\{a\in K:v(a)>0\}\cup\{0\}.
\]
The monomial convention is $t^\gamma=\omega^{-\gamma}$. The field $K$ is algebraically closed \cite{Poonen}. Every finite collection of surcomplex Hahn data can be placed in such a set-sized workspace, and the results below are compatible with enlarging it.

Let $\RR_n=\RR_n(\Gamma)$ be the ring of germs at the ordinary origin of common-domain Hahn families. Its elements are functions on the \emph{whole} monad $\mm_K^n$, not just germs at one fine-topological point. The main results are as follows.

\paragraph{Finite algebra and support-controlled division.}
If $b_1,\ldots,b_\mu$ are polynomial representatives of any basis of the ordinary local algebra of $f$, then their classes are a $K$-basis of
\[
 A_F=\RR_n/(F_1,\ldots,F_n).
\]
For every $G$ one has
\[
 G=\sum_{j=1}^{\mu}b_jN_j(G)+\sum_{i=1}^{n}F_iQ_i(G),
\]
with the bound \eqref{eq:masterbound}. All higher homology of the Hahn Koszul complex vanishes. This is proved in \cref{thm:conservation}.

\paragraph{Actual roots, multiplicities, and base change.}
The common zeros form a finite set in $\mm_K^n$, and the sum of their formal local intersection multiplicities is exactly $\mu$. Enlarging the value group neither changes the multiplication matrices nor creates additional roots. Thus the count holds in the full surcomplex monad as well; see \cref{thm:rootcount,thm:basechange}.

\paragraph{Finite analytic geometry.}
The ring $\RR_n$ is a Noetherian Jacobson domain of dimension $n$. Its maximal ideals are exactly the infinitesimal evaluation ideals, its local rings at these ideals are regular, and their completions are $K[[h_1,\ldots,h_n]]$. Noether normalization yields finite surjective projections of irreducible monad zero sets and a generically reduced finite-fiber statement; see \cref{sec:ring,sec:finite}.

\paragraph{Residue duality through collisions.}
A fixed ordinary residue cycle defines a Hahn functional $\Lambda_F$ on $A_F$. The pairing $(a,b)\mapsto\Lambda_F(ab)$ is perfect, and
\begin{equation}\label{eq:introtrace}
 \Tr(m_G)=\Lambda_F(GJ_F),\qquad
 J_F=\det\left(\frac{\partial F_i}{\partial z_j}\right).
\end{equation}
In particular $\Lambda_F(J_F)=\mu$. For simple roots,
\[
 \Lambda_F(G)=\sum_{F\sh(a)=0}\frac{G\sh(a)}{J_F\sh(a)}.
\]
The perfect pairing, unlike this last simple-root formula, remains valid at a collision. See \cref{sec:residues}.

\subsection{Predecessors and the status of novelty}\label{sec:novelty}
A literature search changes the appropriate framing of these results. Cluckers--Lipshitz--Robinson prove multivariable preparation and division for nonstandard analytic structures \cite[Theorem 2.5 and Remark 2.6]{CLR}. Cluckers--Lipshitz explicitly treat analytic coefficient families with a common radius and arbitrary well-ordered Hahn supports \cite[\S3.1, Remark 3.1.4, Example 3.3(6)]{CL}; their separated analytic theory also contains normalization and Nullstellensatz theorems \cite[\S5.2]{CL}. These are close and substantial predecessors. Passing to complex coefficients, phrasing their consequences in the supplied manuscript's terminology, or giving another Weierstrass recursion is not by itself a novelty claim.

Conservation of intersection number and residue duality are classical in complex analytic geometry \cite{GR,GH,SS}. The underlying perturbation mechanism is a special case of homological perturbation theory \cite{Crainic}. The new-work claim here is deliberately narrower: we give an explicit fixed-neighborhood realization of that mechanism for arbitrary Hahn supports, retain a specified basis throughout the deformation, obtain uniform support certificates for its normal forms and matrices, prove compatibility with all divisible workspace enlargements, and identify the associated moving-root residue functional including its collision-stable pairing.

To the best of the present literature check, this exact combined statement and proof were not found in a published source. This is a proposed original synthesis and formulation, not a certified first theorem in multivariable non-Archimedean analysis. In particular, this article resolves a concrete direction posed in \cite{Source}; it does not claim to settle a separately named, previously recognized open conjecture. The proofs, rather than a priority assertion, are the substantive deliverable.

\section{Hahn support calculus and the common-domain ring}\label{sec:support}
\subsection{Strong summability}
A $K$-element has a unique expression $a=\sum_\gamma a_\gamma t^\gamma$, with $a_\gamma\in\C$ and well-ordered support. For $a\ne0$, $v(a)$ is its least exponent. A family of Hahn series is \emph{strongly summable} when the union of its supports is well ordered and every exponent receives only finitely many nonzero contributions. Its sum is defined coefficientwise.

We use the following standard support facts \cite{Neumann,Poonen}.
\begin{lemma}[Hahn--Neumann support lemma]\label{lem:neumann}
If $S_1,S_2$ are well-ordered subsets of an ordered abelian group, $S_1+S_2$ is well ordered and every exponent has only finitely many representations $s_1+s_2$ with $s_i\in S_i$. If $T\subset\Gamma_{>0}$ is well ordered, $T^*$ is well ordered, and each exponent has finitely many representations as a word in $T$, including its possible word lengths. Consequently, for well-ordered $S$, the coefficients associated with pairs consisting of an element of $S$ and a word in $T$ are locally finite.
\end{lemma}
The empty word represents zero. Positivity of $T$ is essential: zero increments would allow infinitely many contributions at one exponent. These facts apply to supports of vector-valued series as well as scalar-valued series.

\begin{definition}[Hahn families of vectors]
For a complex vector space $V$, let $V((t^\Gamma))$ denote well-ordered Hahn families with coefficients in $V$. A complex-linear map $L:V\to W$ acts coefficientwise and never increases support. This assertion makes no continuity assumption on $L$.
\end{definition}
The last observation will permit the use of ordinary holomorphic division maps on a fixed vector space $\OO(U)$ even when the coefficient functions have arbitrarily large ordinary norms.

\begin{lemma}[Positive operator inversion]\label{lem:operator}
Let $T\subset\Gamma_{>0}$ be well ordered, and let
\[
 \Delta=\sum_{\gamma\in T}t^\gamma L_\gamma
\]
act on $V((t^\Gamma))$, where $L_\gamma:V\to V$ is complex-linear. Then
\begin{equation}\label{eq:neumannoperator}
 (1+\Delta)^{-1}=\sum_{k\ge0}(-\Delta)^k
\end{equation}
is a well-defined $K$-linear operator. Applied to a vector of support $S$, its output has support contained in $S+T^*$.
\end{lemma}
\begin{proof}
At an exponent $\nu$, an expression in the right side uses an input coefficient at $s\in S$ and a word $(\gamma_1,\ldots,\gamma_k)$ satisfying $s+\gamma_1+\cdots+\gamma_k=\nu$. There are finitely many such choices by \cref{lem:neumann}. Thus the coefficient is a finite sum of compositions of the $L_\gamma$. The total support is well ordered. The usual cancellation in $(1+\Delta)\sum_{k\ge0}(-\Delta)^k$ is legitimate coefficient by coefficient and gives the identity. The same proof works on the other side. Commuting an additional scalar Hahn series through these operators is justified by the same finite-contribution property, which proves $K$-linearity.
\end{proof}

\begin{example}[Why this is not a convergence proof]\label{ex:ranktwo}
Order $\Q\oplus\Q$ lexicographically. For $e=(0,1)$ and $b=(1,0)$, every $ne$ is smaller than $b$. Nevertheless $\sum_{n\ge0}t^{ne}$ is a legitimate Hahn series and is $(1-t^e)^{-1}$. Its partial sums do not approach this value in the valuation topology: their errors never have valuation greater than $b$. Thus replacing \cref{lem:operator} by an assertion that the errors tend to zero would lose legitimate cases of the theorem.
\end{example}

\subsection{Coherent functions in several variables}
For an ordinary connected open set $U\subset\C^n$, define
\[
 \HH_\Gamma(U)=\left\{\sum_{\gamma\in S}f_\gamma t^\gamma:
 f_\gamma\in\OO(U),\ S\subset\Gamma\text{ well ordered}\right\}.
\]
Zero coefficient functions may be omitted. Let
\[
 U\sh_K=\{c+\eps:c\in U,\ \eps\in\mm_K^n\}.
\]
The evaluation of $F=\sum f_\gamma t^\gamma$ is
\begin{equation}\label{eq:evaluation}
 F\sh(c+\eps)=
 \sum_{\substack{\gamma\in\supp F\\\alpha\in\N^n}}
 t^\gamma\frac{\partial^\alpha f_\gamma(c)}{\alpha!}\eps^\alpha.
\end{equation}
Here $\alpha$ is a multiindex, $\alpha!=\prod_j\alpha_j!$, and $\eps^\alpha=\prod_j\eps_j^{\alpha_j}$.

\begin{proposition}[Evaluation and support]\label{prop:evaluation}
Formula \eqref{eq:evaluation} is strongly summable. Evaluation respects sums, products, and coefficientwise partial derivatives. It is injective on $\HH_\Gamma(U)$. Composition with an infinitesimal translation or with a positive-support perturbation of an ordinary holomorphic map is well defined whenever the ordinary domains match.
\end{proposition}
\begin{proof}
Put $P=\bigcup_j\supp\eps_j$. All exponents in the evaluation lie in $\supp F+P^*$. The multiindices of each fixed total degree give only finitely many choices, so \cref{lem:neumann} proves local finiteness. The ordinary Taylor product and composition identities may therefore be rearranged coefficientwise. Differentiation of coefficient functions does not introduce new exponents. At ordinary points, $F\sh(c)=\sum f_\gamma(c)t^\gamma$; if all these values vanish, uniqueness of Hahn coefficients makes every $f_\gamma$ identically zero. The same support argument proves the composition assertion.
\end{proof}
As in the supplied manuscript, these are genuine fine-analytic evaluations. The results below are principally statements in the coherent coefficient algebra, so they do not need a fine-topological compactness or path theorem.

\begin{definition}[Common-domain monad germ ring]\label{def:ring}
Let
\begin{equation}\label{eq:ring}
 \RR_n(\Gamma)=\varinjlim_{U\ni0}\HH_\Gamma(U),
\end{equation}
where the limit is over connected ordinary neighborhoods of $0\in\C^n$ and the maps are restrictions. We set $\RR_0(\Gamma)=K$.
\end{definition}
One may use only ordinary polydiscs in \eqref{eq:ring}. Restriction is injective by the ordinary identity theorem, and multiplication of leading coefficient germs shows that $\RR_n$ is a domain. Every element evaluates on all of $\mm_K^n$.

The quantifier order in \eqref{eq:ring} matters. All coefficient functions of one representative have \emph{one} common domain. The larger formal expression $\OO_{\C^n,0}((t^\Gamma))$ would permit independently shrinking coefficient domains and is not the ring used here. For example,
\[
 \sum_{m\ge1}\frac{t^m}{1-mz}
\]
has a holomorphic germ at zero at each exponent, but there is no positive ordinary radius on which all its displayed coefficient functions are holomorphic. It is not an element of $\RR_1$ with those coefficients.

\begin{lemma}[Infinitesimal translation and coordinate quotients]\label{lem:translation}
For $a\in\mm_K^n$, translation $T_aF(z)=F(z+a)$ is an automorphism of $\RR_n$ with inverse $T_{-a}$. Moreover
\begin{equation}\label{eq:coordinatequotient}
 \RR_n/(z_1-a_1,\ldots,z_k-a_k)\simeq\RR_{n-k}
 \quad(0\le k\le n).
\end{equation}
In particular evaluation at $a$ is surjective onto $K$, with kernel
$\mm_a=(z_1-a_1,\ldots,z_n-a_n)$.
\end{lemma}
\begin{proof}
The translation is formed by the Taylor formula with support bounded by the original support plus the monoid of the coordinates of $a$. Each derivative coefficient remains holomorphic on the same ordinary domain. The binomial identity proves $T_aT_{-a}=1$. For $a=0$, the ordinary identity
\[
 h(z_1,z')=h(0,z')+z_1\frac{h(z_1,z')-h(0,z')}{z_1}
\]
has a holomorphic quotient on any polydisc. Applying it coefficientwise eliminates one coordinate without changing support; iterating eliminates $k$. Translation gives the statement for general $a$.
\end{proof}

\section{Preparation and division with holomorphic parameters}\label{sec:prep}
This section supplies the common-domain version needed later. It is a complex-coefficient presentation of a Weierstrass mechanism with established nonstandard-analytic predecessors \cite{CLR,CL}, and a parameter-dependent extension of the recursion in \cite{Source}.

\subsection{One ordinary division pair on one fixed domain}
Write $z=(x,w)$ with $x\in\C^{n-1}$. Let $V$ be an ordinary polydisc and let $D_R=\{|w|<R\}$. Suppose $P_0(x,w)$ is monic of degree $d$ in $w$, with coefficients holomorphic on $V$, and every root of $P_0(x,\cdot)$ lies in $D_r$, where $r<R$ is fixed.

Choose $r<\rho<R$. For $h\in\OO(V\times D_R)$, put
\begin{equation}\label{eq:ordinaryR}
 R_0h(x,w)=\frac1{2\pi i}\int_{|\zeta|=\rho}
 \frac{h(x,\zeta)}{P_0(x,\zeta)}
 \frac{P_0(x,\zeta)-P_0(x,w)}{\zeta-w}\,d\zeta.
\end{equation}
The quotient involving $P_0$ is a polynomial in $w$ of degree less than $d$. Consequently $R_0h$ has that degree, with coefficients holomorphic on $V$. Define
\begin{equation}\label{eq:ordinaryD}
 D_0h=\frac{h-R_0h}{P_0}.
\end{equation}
For $|w|<\rho$, the Cauchy formula shows that $h-R_0h$ is divisible by $P_0$. All zeros of $P_0$ are in this smaller disc, so the quotient extends holomorphically to all of $V\times D_R$. We obtain
\begin{equation}\label{eq:ordinarydivision}
 h=R_0h+P_0D_0h,\qquad \deg_wR_0h<d.
\end{equation}
Uniqueness follows because a polynomial of degree less than $d$ divisible by the monic degree-$d$ polynomial $P_0(x,\cdot)$ must vanish. Both operators are complex-linear and act on \emph{the same} space $\OO(V\times D_R)$.

\subsection{Support-controlled preparation}
\begin{theorem}[Parameterized Hahn preparation]\label{thm:prep}
Let $F\in\RR_n$ be nonzero. Normalize its least Hahn exponent so that
\[
 F=t^\lambda(f_0+E),\qquad v_H(E)>0,
\]
allowing $E=0$. Suppose $f_0(0,w)$ has a zero of order $d\ge1$ at $w=0$. After choosing an ordinary product neighborhood, ordinary preparation gives $f_0=u_0P_0$, where $u_0$ is nowhere zero and $P_0$ is monic of degree $d$, with $P_0(0,w)=w^d$. On a possibly smaller product neighborhood, there are unique normalized factors
\begin{equation}\label{eq:prep}
 F=t^\lambda u_0(P_0+A)(1+B),\qquad \deg_w A<d,
\end{equation}
where $A$ and $B$ have strictly positive Hahn support. If $T$ is the support of $E/u_0$, their supports are contained in $T^*\setminus\{0\}$. Every coefficient function of every factor is holomorphic on this one neighborhood.
\end{theorem}
\begin{proof}
Ordinary Weierstrass preparation is used only once, on $f_0$ \cite{GR}. Shrink the parameter polydisc so the roots of $P_0$ stay in a fixed smaller $w$-disc, giving the operators \eqref{eq:ordinaryR}--\eqref{eq:ordinaryD}. Divide by $t^\lambda u_0$ and write the normalized error as $E'$. The equation becomes
\[
 E'=A+P_0B+AB.
\]
Proceed by well-founded recursion on $T^*\setminus\{0\}$. At exponent $\nu$, define
\begin{align}
 h_\nu&=E'_\nu-
   \sum_{\substack{\alpha+\beta=\nu\\\alpha,\beta>0}}
        A_\alpha B_\beta,\label{eq:preprecursion1}\\
 A_\nu&=R_0h_\nu,\qquad B_\nu=D_0h_\nu.\label{eq:preprecursion2}
\end{align}
The sum is finite and uses only smaller exponents. The division identity proves the desired equation at $\nu$. The support lemma proves that the assembled factors exist. Crucially, $R_0,D_0$ never change the chosen ordinary domain.

For uniqueness, compare two normalized pairs and take the least exponent where either differs. The product $AB$ cannot contribute a first difference at that exponent, because both factors have positive support. Thus $\Delta A_\nu+P_0\Delta B_\nu=0$. Ordinary division uniqueness makes both differences zero, a contradiction. The factor $1+B$ is a unit by positive-support geometric inversion, and $u_0$ is an ordinary unit.
\end{proof}

\begin{theorem}[Parameterized Hahn division]\label{thm:division}
Let $P=P_0+A$ be the monic factor in \cref{thm:prep}, and let $G$ be a Hahn-holomorphic datum on its product neighborhood. There are unique $Q$ and $R$ on that same neighborhood with
\begin{equation}\label{eq:division}
 G=PQ+R,\qquad \deg_wR<d.
\end{equation}
Writing $T_A=\supp A\subset\Gamma_{>0}$, both outputs have support contained in $\supp G+T_A^*$. They are given by
\begin{equation}\label{eq:divisionformula}
 Q=\sum_{k\ge0}\bigl(-D_0(A\,\cdot)\bigr)^kD_0G,
 \qquad R=R_0(G-AQ).
\end{equation}
For an arbitrary germ $G\in\RR_n$, one first chooses a common smaller product neighborhood; no subsequent shrinking is needed.
\end{theorem}
\begin{proof}
Apply $D_0$ to \eqref{eq:division}. The equation for $Q$ is
\[
 Q+D_0(AQ)=D_0G.
\]
The positive operator inversion lemma gives \eqref{eq:divisionformula} and its support bound. Set $R=R_0(G-AQ)$ and use ordinary division to recover \eqref{eq:division}. If $PQ+R=0$ with $\deg_wR<d$, the same invertible operator gives $Q=0$ and then $R=0$.
\end{proof}

\begin{corollary}[A finite free hypersurface quotient]\label{cor:freehypersurface}
For $P$ as above,
\[
 \RR_n/(P)\simeq
 \RR_{n-1}\oplus\RR_{n-1}w\oplus\cdots\oplus\RR_{n-1}w^{d-1}
\]
as an $\RR_{n-1}$-module. For each infinitesimal parameter $a\in\mm_K^{n-1}$, all $d$ roots of $P\sh(a,w)$ are infinitesimal, counted with multiplicity.
\end{corollary}
\begin{proof}
Division gives the displayed module isomorphism. The lower coefficients of $P_0(a,w)$ are infinitesimal because their ordinary values at $a=0$ vanish; the coefficients of $A(a,w)$ are also infinitesimal. If a root $w$ were not infinitesimal, divide the equation by $w^d$. The result would be $1$ plus an infinitesimal equal to zero. Algebraic closedness of $K$ supplies all $d$ roots.
\end{proof}

\section{The algebraic geometry of a monad}\label{sec:ring}
The results of this section are structural adaptations of the Weierstrass-algebra method. Their role is to ensure that the normal forms constructed later describe \emph{all actual points}, not just an abstract quotient. Analytic normalization and Nullstellensatz theorems in related non-Archimedean systems are established in \cite[\S5.2]{CL}.

\subsection{Noetherianity and dimension}
\begin{theorem}[Noetherian monad ring]\label{thm:noetherian}
The ring $\RR_n$ is a Noetherian domain of Krull dimension $n$.
\end{theorem}
\begin{proof}
The domain assertion follows from multiplication of leading coefficient germs. For Noetherianity, use induction on $n$, starting with $\RR_0=K$. Let $I\subset\RR_n$ be a nonzero ideal and choose $0\ne F\in I$. An ordinary invertible linear change of coordinates makes its leading ordinary coefficient regular in the last variable: choose a direction where its first nonzero homogeneous Taylor polynomial does not vanish. If this coefficient has order zero, $F$ is a unit by geometric inversion after shrinking the ordinary neighborhood, and $I=\RR_n$.

Otherwise preparation replaces $F$, up to a unit, by a monic polynomial $P$ of positive degree in the last variable. By \cref{cor:freehypersurface}, $\RR_n/(F)$ is a finite free module over $\RR_{n-1}$. Since the latter ring is Noetherian, its submodule $I/(F)$ is finitely generated. Lifting those generators and adjoining $F$ generates $I$ as an $\RR_n$-ideal.

For dimension, the coordinate ideals
\[
 (0)\subsetneq(z_1)\subsetneq(z_1,z_2)\subsetneq\cdots
 \subsetneq(z_1,\ldots,z_n)
\]
are prime by \cref{lem:translation}, giving the lower bound $n$. For the upper bound, prepend $(0)$ to any prime chain when necessary and choose a nonzero element in its first nonzero prime. Preparation as above makes the quotient by this element finite integral over $\RR_{n-1}$. Integral extensions preserve dimension; therefore the remaining chain has length at most $n-1$. This gives dimension at most $n$. The elementary integral-extension facts used here, including incomparability, going up, and lying over, are the standard commutative-algebra results recorded in \cite{Stacks}.
\end{proof}

\begin{proposition}[Noether normalization]\label{prop:normalization}
If $\pp$ is a prime of $\RR_n$ and $d=\dim(\RR_n/\pp)$, an ordinary linear coordinate system can be chosen so that the first $d$ coordinates induce a finite injective homomorphism
\begin{equation}\label{eq:normalization}
 \RR_d\hookrightarrow A=\RR_n/\pp.
\end{equation}
\end{proposition}
\begin{proof}
Induct on $n$. When $\pp=0$, take $d=n$. Otherwise prepare a nonzero element of $\pp$ in the last variable. Then $A$ is finite over $\RR_{n-1}/\mathfrak q$, where $\mathfrak q$ is the contraction of $\pp$. By induction, this domain is finite over an injective coordinate copy of some $\RR_d$. Finiteness is transitive. All coordinate changes in the induction act ordinarily on the remaining coordinates, so their composition is an ordinary linear change on $\C^n$. Dimension is preserved in each finite integral extension, which identifies the final number of coordinates with $\dim A$.
\end{proof}

\subsection{Every maximal ideal is an actual infinitesimal point}
\begin{theorem}[Weak monad Nullstellensatz]\label{thm:weaknull}
Every maximal ideal of $\RR_n$ is uniquely of the form
\[
 \mm_a=(z_1-a_1,\ldots,z_n-a_n),\qquad a\in\mm_K^n.
\]
Its residue field is $K$, and its residue map is Hahn evaluation at $a$.
\end{theorem}
\begin{proof}
Induct on $n$. The assertion for $n=0$ is immediate. For $n>0$, $\RR_n$ is not a field, since evaluation at zero shows that $z_1$ is a nonunit. Thus a maximal ideal $\mathfrak m$ is nonzero. Prepare an element of $\mathfrak m$ and use the resulting finite extension from $\RR_{n-1}$. The contraction $\mathfrak n$ is maximal: a subring over which a field is integral is itself a field. By induction $\RR_{n-1}/\mathfrak n=K$. The field $\RR_n/\mathfrak m$ is a finite extension of $K$, hence is $K$ because $K$ is algebraically closed.

Let $a_j\in K$ be the image of $z_j$. If $v(a_j)<0$, then $z_j-a_j$ has a nonzero constant as its leading Hahn coefficient after normalization, and is a unit. If $v(a_j)=0$, its leading coefficient is $z_j-c$ for some $c\in\C^\times$, which is a unit on a sufficiently small ordinary neighborhood of zero. Both cases contradict $z_j-a_j\in\mathfrak m$. Therefore every $a_j$ is infinitesimal or zero. By \cref{lem:translation}, the ideal generated by these coordinate differences already has quotient $K$ and is precisely the evaluation kernel. It must equal $\mathfrak m$. The coordinate images give uniqueness.
\end{proof}

\begin{lemma}[Infinitesimal evaluation detects nonzero germs]\label{lem:density}
For every nonzero $F\in\RR_n$, some $a\in\mm_K^n$ satisfies $F\sh(a)\ne0$.
\end{lemma}
\begin{proof}
Normalize $F=t^\lambda(f_0+E)$ with $f_0\ne0$ and positive-support $E$. Let $m$ be the least total degree of a nonzero homogeneous term of $f_0$. Choose $c\in\C^n$ at which that homogeneous polynomial is nonzero. If $E\ne0$, put $e=\min\supp E>0$ and choose $\delta>0$ in $\Gamma$ with $m\delta<e$, for example $\delta=e/(m+1)$. If $E=0$, any positive $\delta$ suffices. Such a $\delta$ exists because $\Gamma$ is nonzero and divisible.

At $a=t^\delta c$, the leading term of $f_0\sh(a)$ has valuation $m\delta$. Every value contributed by $E$ has valuation at least $e$. It cannot cancel that term. This proves $F\sh(a)\ne0$.
\end{proof}
In particular the map from $\RR_n$ to functions on the whole monad is injective. The nontriviality assumption on $\Gamma$ is necessary here: when $\Gamma=0$, the monad has only its origin and cannot detect ordinary analytic germs.

\begin{theorem}[Strong monad Nullstellensatz]\label{thm:strongnull}
For an ideal $I\subset\RR_n$, put
\[
 V(I)=\{a\in\mm_K^n:F\sh(a)=0\text{ for every }F\in I\}.
\]
Then
\begin{equation}\label{eq:strongnull}
 I(V(I))=\sqrt I.
\end{equation}
Equivalently, $\RR_n$ is Jacobson and all its closed points are its infinitesimal evaluation points.
\end{theorem}
\begin{proof}
We first show that a prime $\pp$ is the intersection of the maximal ideals containing it. Let $A=\RR_n/\pp$ and let $0\ne g\in A$. Normalize $A$ over $B=\RR_d$ as in \cref{prop:normalization}. The finite-dimensional domain
\[
 A\otimes_B\Frac(B)
\]
is a field. Consequently $g^{-1}$ belongs to it. Clearing a denominator gives an equation $gh=b$ with $h\in A$ and $0\ne b\in B$.

By \cref{lem:density}, choose an evaluation maximal ideal $\mathfrak n$ of $B$ not containing $b$. Lying over for the integral extension $B\subset A$ gives a maximal ideal $M$ of $A$ over $\mathfrak n$. The equation $gh=b$ implies $g\notin M$. Pulling $M$ back to $\RR_n$ and applying \cref{thm:weaknull} produces a point of $V(\pp)$ where any representative of $g$ is nonzero. This proves the prime-intersection assertion.

If $F\notin\sqrt I$, choose a prime $\pp\supset I$ not containing $F$. The preceding argument finds an evaluation maximal ideal containing $\pp$ but not $F$. Thus $F$ does not vanish on all of $V(I)$. The reverse inclusion in \eqref{eq:strongnull} is immediate in the field $K$.
\end{proof}
This is a statement about a ring of functions on an entire infinitesimal monad. It is not the ordinary complex local Nullstellensatz restated with a different scalar field: $\RR_n$ has many maximal ideals, whereas $\C\{z\}$ is local.

\subsection{Local regularity and formal multiplicity}
\begin{proposition}[Local coordinates and completion]\label{prop:completion}
For $a\in\mm_K^n$, the local ring $(\RR_n)_{\mm_a}$ is regular of dimension $n$, and
\begin{equation}\label{eq:completion}
 \widehat{(\RR_n)_{\mm_a}}\simeq K[[h_1,\ldots,h_n]],
 \qquad h_j=z_j-a_j.
\end{equation}
The image of $F$ in this completion is its formal Taylor series at $a$.
\end{proposition}
\begin{proof}
Translate $a$ to zero. The maximal ideal is generated by the $n$ coordinate functions. The coordinate prime chain survives localization, so the local dimension is at least $n$ and hence exactly $n$ by \cref{thm:noetherian}. A Noetherian local ring of dimension $n$ whose maximal ideal has $n$ generators is regular.

For every integer $r\ge1$, ordinary Taylor division coefficientwise gives a unique decomposition into the polynomial of total degree less than $r$ and a member of $(z_1,\ldots,z_n)^r$. The coefficients of the finite Taylor polynomial are elements of $K$. This proves
\[
 \RR_n/\mm_0^r\simeq K[z_1,\ldots,z_n]/(z_1,\ldots,z_n)^r.
\]
Localizing does not change this quotient, whose radical is already maximal. Taking the inverse limit over $r$ proves \eqref{eq:completion}. Translation and \cref{prop:evaluation} identify its coefficients with the Taylor coefficients at $a$.
\end{proof}
For an isolated zero of $F_1,\ldots,F_n$ at $a$, the local intersection multiplicity can consequently be defined without ambiguity as
\begin{equation}\label{eq:localmultiplicity}
 \mu_a(F)=\dim_K
 \frac{K[[h_1,\ldots,h_n]]}
      {(F_1(a+h),\ldots,F_n(a+h))},
\end{equation}
when finite. It equals the length of the corresponding local quotient of $\RR_n$, since completion preserves finite-length modules.

\section{Finite maps of monad zero sets}\label{sec:finite}
\begin{theorem}[Finite projection and generic fibers]\label{thm:finitemap}
Let $\pp\subset\RR_n$ be prime and $A=\RR_n/\pp$ have dimension $d$. After an ordinary linear coordinate change, the coordinate projection
\[
 \pi:V(\pp)\longrightarrow\mm_K^d
\]
is surjective and has finite fibers. It is induced by a finite injective map $B=\RR_d\hookrightarrow A$. If $A$ is generated by $N$ elements as a $B$-module, every fiber has scheme length at most $N$.

Let $r=[\Frac(A):\Frac(B)]$. There is a nonzero $s\in B$ such that every $b\in\mm_K^d$ with $s\sh(b)\ne0$ has exactly $r$ distinct preimages, each reduced. Such $b$ exist.
\end{theorem}
\begin{proof}
Use \cref{prop:normalization}. Lying over makes the map on maximal ideals surjective. By \cref{thm:weaknull}, these are exactly the evaluation points on both sides, proving the asserted surjectivity on monads. For an evaluation maximal ideal $\mathfrak n_b\subset B$, the fiber algebra $A/\mathfrak n_bA$ is a nonzero finite-dimensional $K$-algebra. The $N$ module generators span it as a vector space, giving the bound and finiteness of the fiber.

Set $L=\Frac(A)$ and $L_0=\Frac(B)$. The extension $L/L_0$ is finite and separable, since the characteristic is zero. By the primitive element theorem and clearing denominators, choose a primitive element $\theta$ in a localization of $A$. The finitely many module generators of $A$ can be written as $L_0$-linear combinations of $1,\theta,\ldots,\theta^{r-1}$. Localize at one nonzero $s\in B$ so that these expressions, the coefficients of the monic minimal polynomial $q(T)$ of $\theta$, and the element $\theta$ itself are all defined over $B_s$. Enlarge $s$ to make the nonzero discriminant of $q$ invertible. Then
\[
 A_s\simeq B_s[T]/(q),\qquad \deg q=r,\qquad\disc(q)\in B_s^\times.
\]
Indeed, the expressions for the generators give surjectivity; monic division and independence of $1,\theta,\ldots,\theta^{r-1}$ over $L_0$ give injectivity. Each stated fiber is therefore the algebra of a degree-$r$ polynomial with nonzero discriminant over the algebraically closed field $K$, and has $r$ distinct points. Finally \cref{lem:density} supplies a $b$ avoiding $s=0$.
\end{proof}
The word \emph{finite} here has its algebraic meaning: finite module and finite fibers. No claim of properness, compactness, or classical path lifting in the full surreal fine topology is included. This is the finite-mapping statement naturally attached to the coherent ring of the supplied framework.

\section{A support-controlled conjugacy of complexes}\label{sec:contraction}
We now prove the mechanism behind conservation of multiplicity. It is a particularly transparent specialization of the homological perturbation lemma \cite{Crainic}; the explicit proof records the sign convention and the exact support control.

\subsection{An ordinary contraction}
Let $(C_\bullet,d)$ be a complex of complex vector spaces, concentrated in nonnegative degrees, with homology only in degree zero. Let $V=H_0(C_\bullet)$ be finite-dimensional. Choose linear maps
\[
 i:V\to C_0,\qquad p:C_0\to V,\qquad h:C_k\to C_{k+1}
\]
and extend $p$ by zero in other degrees, so that
\begin{equation}\label{eq:contraction}
 pi=1,\quad di=0,\quad pd=0,\quad
 dh+hd=1-ip,\quad h^2=0,\quad hi=0,\quad ph=0.
\end{equation}
Such data always exist. To see this, split each space into boundaries and a complement mapping isomorphically onto the boundaries one degree lower; in degree zero also split off the selected representatives $i(V)$. Define $h$ to invert $d$ on boundaries and to vanish on the chosen complements and on $i(V)$. This construction can enforce any specified representatives of a basis of $V$.

These are vector-space splittings. If $C_k$ consists of holomorphic functions on a domain, no boundedness or continuity of these splittings is asserted. Their outputs nevertheless belong to the same specified space of holomorphic functions.

\subsection{The perturbation identity}
\begin{theorem}[Positive Hahn conjugacy]\label{thm:conjugacy}
Extend the data \eqref{eq:contraction} coefficientwise to Hahn families. Let
\[
 \delta=\sum_{\gamma\in T}t^\gamma\delta_\gamma
\]
be a degree-$-1$ operator, with $T\subset\Gamma_{>0}$ well ordered, and suppose $(d+\delta)^2=0$. Set $D=d+\delta$ and
\begin{equation}\label{eq:perturbedmaps}
 U=1+\delta h,\qquad
 p_F=pU^{-1},\qquad h_F=hU^{-1}.
\end{equation}
Then
\begin{equation}\label{eq:chainconjugacy}
 D=UdU^{-1},\qquad Ui=i,
\end{equation}
and
\begin{equation}\label{eq:perturbedcontraction}
 p_Fi=1,\qquad p_FD=0,\qquad
 Dh_F+h_FD=1-ip_F.
\end{equation}
All these operators are well defined. Applied to input support $S$, $U^{-1}$, $p_F$, and $h_F$ have output support contained in $S+T^*$.
\end{theorem}
\begin{proof}
The inverse of $U$ is \cref{lem:operator}. Because the complex is concentrated in nonnegative degrees and $i(V)$ lies in degree zero, $\delta i=0$. From $D^2=0$ we have $d\delta+\delta d+\delta^2=0$. A direct calculation gives
\begin{align*}
 Ud-DU
 &=\delta hd-\delta-d\delta h-\delta^2h\\
 &=\delta hd-\delta+\delta dh\\
 &=\delta(hd+dh-1)=-\delta ip=0.
\end{align*}
Thus $Ud=DU$. Moreover $Ui=i$ because $hi=0$, proving \eqref{eq:chainconjugacy}. Since $h^2=0$, $Uh=h$, and hence $UhU^{-1}=h_F$. Conjugating $dh+hd=1-ip$ by $U$ gives the last identity in \eqref{eq:perturbedcontraction}. The other two follow from $U^{-1}i=i$ and $U^{-1}D=dU^{-1}$. Every identity is an equality of strongly summable coefficients, justified by \cref{lem:neumann}. The support bound follows from the geometric inverse for $U$ and from the fact that $p,h$ preserve support.
\end{proof}
The absence of a transferred differential on $V$ is not an extra assumption: it follows from concentration in degree zero. In a complex with homology in several degrees, positive perturbations can produce a nontrivial differential on the homology space, so the theorem would have to be modified.

\section{Conservation of an isolated complete intersection}\label{sec:conservation}
\subsection{The ordinary fixed-domain input}
Let $f_1,\ldots,f_n$ be ordinary holomorphic germs with an isolated common zero at zero. Their local algebra
\[
 A_0=\C\{z_1,\ldots,z_n\}/(f_1,\ldots,f_n)
\]
has finite dimension $\mu$. Choose polynomial representatives $b_1,\ldots,b_\mu$ of a basis. Such representatives exist because a power of the maximal ideal is contained in $(f_1,\ldots,f_n)$.

Choose a sufficiently small ordinary Stein polydisc $\Omega$ on which all $f_i$ are holomorphic and have no common zero except zero. Consider the global-section Koszul complex
\begin{equation}\label{eq:koszul}
 C_k=\OO(\Omega)\otimes_\C\bigwedge^k\C^n,
\end{equation}
with differential
\[
 d_f(g\,e_{i_1}\wedge\cdots\wedge e_{i_k})
 =\sum_{j=1}^{k}(-1)^{j-1}f_{i_j}g\,
   e_{i_1}\wedge\cdots\widehat e_{i_j}\cdots\wedge e_{i_k}.
\]

\begin{lemma}[Ordinary contraction on a fixed neighborhood]\label{lem:ordinarykoszul}
The complex \eqref{eq:koszul} has no positive-degree homology and its degree-zero homology is $A_0$. It therefore admits a contraction \eqref{eq:contraction} with $V=\C^\mu$ and
\[
 i(c_1,\ldots,c_\mu)=\sum_jc_jb_j.
\]
Every component of its maps takes holomorphic functions on $\Omega$ to holomorphic functions on that same $\Omega$.
\end{lemma}
\begin{proof}
At zero, the $n$ germs generate an ideal of finite colength in the regular local ring of dimension $n$. They are a system of parameters and a regular sequence, so their Koszul complex is exact in positive degrees. At any other point of $\Omega$, one $f_i$ is a unit in the stalk and the Koszul complex is contractible. Thus the sheaf Koszul complex resolves a coherent quotient sheaf supported just at zero, whose stalk there is $A_0$.

Kernels in this finite complex are coherent. Cartan's theorem B on the Stein domain $\Omega$ makes their higher cohomology vanish; applying it to the successive short exact sequences shows that global sections preserve the required exactness. The global sections of the quotient sheaf are precisely $A_0$. These ordinary analytic and regular-local facts are classical \cite{GR,Stacks}. The vector-space splitting described in \cref{sec:contraction} now provides the asserted maps. Their codomains are the fixed spaces \eqref{eq:koszul}, which is the required common-domain statement.
\end{proof}
This use of Cartan B is substantive. Choosing a separate germ-level division for each coefficient and then intersecting all its neighborhoods would not prove the lemma's final assertion.

\subsection{The normal-form theorem}
\begin{theorem}[Support-explicit conservation of multiplicity]\label{thm:conservation}
Let $f$, $\mu$, and $b_1,\ldots,b_\mu$ be as above. Suppose $F_i=f_i+E_i$ with all $E_i\in\RR_n$ of strictly positive Hahn support, allowing zero errors, and let
\[
 T=\bigcup_{i=1}^{n}\supp E_i\subset\Gamma_{>0}.
\]
Then the following statements hold.
\begin{enumerate}
\item The classes of $b_1,\ldots,b_\mu$ form a $K$-basis of
$A_F=\RR_n/(F_1,\ldots,F_n)$. In particular $\dim_KA_F=\mu$.
\item For $G\in\RR_n$ there are $K$-linear normal-form coordinates $N_j(G)$ and division coefficients $Q_i(G)\in\RR_n$ satisfying
\begin{equation}\label{eq:normalform}
 G=\sum_{j=1}^{\mu}b_jN_j(G)+\sum_{i=1}^{n}F_iQ_i(G).
\end{equation}
On one fixed ordinary domain representing the inputs, these outputs have support contained in $\supp G+T^*$.
\item The higher homology of the Koszul complex of $F_1,\ldots,F_n$ over $\RR_n$ vanishes.
\item The multiplication matrices $M_\ell$ of $z_\ell$ in the prescribed basis have entries supported in $T^*$. Their coefficients at exponent zero are the ordinary multiplication matrices of $z_\ell$ on $A_0$.
\end{enumerate}
The normal-form coordinates are intrinsic once the basis is fixed. The division coefficients may depend on the ordinary contraction chosen in the proof.
\end{theorem}
\begin{proof}
Take one small Stein polydisc $\Omega$ representing all errors, with the ordinary zero set of $f$ equal to $\{0\}$ there. Choose the contraction in \cref{lem:ordinarykoszul}, and extend it coefficientwise to the Hahn Koszul complex over $\HH_\Gamma(\Omega)$. Its ordinary differential is $d_f$. The perturbation
\[
 \delta=d_F-d_f
\]
is the Koszul differential formed from $E_1,\ldots,E_n$. It has support in $T$; each coefficient is a linear operator on the same holomorphic function spaces. Koszul differentials square to zero, so \cref{thm:conjugacy} applies.

In degree zero, define
\begin{equation}\label{eq:normalformoperator}
 N(G)=p(1+\delta h)^{-1}G,
 \qquad Q(G)=h(1+\delta h)^{-1}G\in C_1((t^\Gamma)).
\end{equation}
The contraction identity gives $G=iN(G)+d_FQ(G)$, which is \eqref{eq:normalform}. Also $N\circ i=1$ and $N\circ d_F=0$. Therefore $N$ identifies the quotient with $K^\mu$, with basis the classes of the $b_j$. The conjugacy of complexes proves the positive-degree exactness. Every operation is on the one domain $\Omega$, and \cref{thm:conjugacy} gives the support bound.

For an additional germ $G$, choose a smaller common Stein polydisc if needed and repeat the construction. Restrictions of the quotient algebras send each $b_j$ to the same $b_j$ and are isomorphisms because these are bases on both domains. The normal forms therefore agree under restriction. Passing to the direct limit proves the assertions over $\RR_n$ and shows their independence from the chosen contraction. Positive-degree exactness likewise passes to the limit.

Finally, the $j$th column of $M_\ell$ is $N(z_\ell b_j)$. The input is ordinary, so its support is $\{0\}$, and the output support is in $T^*$. At exponent zero, the inverse in \eqref{eq:normalformoperator} is the identity; hence this column reduces to $p(z_\ell b_j)$, its ordinary normal form.
\end{proof}

\subsection{A coefficient recursion rather than an iteration limit}
Formula \eqref{eq:normalformoperator} has a useful explicit reading. On degree zero let
\[
 L=\delta h=\sum_{\gamma\in T}t^\gamma L_\gamma,
 \qquad W=(1+L)^{-1}G.
\]
Its coefficients are determined by
\begin{equation}\label{eq:Wrecursion}
 W_\nu=G_\nu-
 \sum_{\substack{\gamma\in T,\ \eta\in\supp G+T^*\\
                  \gamma+\eta=\nu}}
 L_\gamma W_\eta.
\end{equation}
All $\eta$ on the right are smaller than $\nu$. For any fixed $\nu$, finitely many input coefficients and finitely many words in perturbation coefficients are involved. Then $N_\nu=pW_\nu$ and $Q_\nu=hW_\nu$.

This is an effective recursion \emph{relative to available ordinary division maps and an effective presentation of the support}. Arbitrary complex-linear splittings and arbitrary well-ordered sets need not come with computable presentations. The theorem does not assert a finite-time enumeration algorithm for every set-theoretically possible Hahn input. The finite-dependence statement at a specified coefficient, however, is unconditional.

\begin{corollary}[No loss of infinitesimal solutions]\label{cor:matrices}
The matrices $M_1,\ldots,M_n$ commute. For each $\ell$,
\begin{equation}\label{eq:characteristicreduction}
 \det(XI-M_\ell)\equiv X^\mu\pmod{\mm_K[X]}.
\end{equation}
Thus every coordinate eigenvalue of the deformed finite algebra is infinitesimal.
\end{corollary}
\begin{proof}
The matrices represent multiplication in a commutative algebra, so they commute. In the ordinary local algebra $A_0$, every coordinate belongs to its nilpotent maximal ideal. Its multiplication matrix is nilpotent, with characteristic polynomial $X^\mu$. The reduction statement follows from \cref{thm:conservation}. A root of a monic polynomial whose lower coefficients are infinitesimal must be infinitesimal by the argument in \cref{cor:freehypersurface}.
\end{proof}

\section{Root counts and independence of the workspace}\label{sec:roots}
\begin{theorem}[Exact multivariable infinitesimal root count]\label{thm:rootcount}
Under the assumptions of \cref{thm:conservation}, the common-zero set
\[
 Z_F=\{a\in\mm_K^n:F_1\sh(a)=\cdots=F_n\sh(a)=0\}
\]
is finite and nonempty. With local multiplicity defined by \eqref{eq:localmultiplicity},
\begin{equation}\label{eq:rootcount}
 \sum_{a\in Z_F}\mu_a(F)=\mu.
\end{equation}
The points are the joint characters of the commuting multiplication matrices $M_1,\ldots,M_n$; equivalently, their coordinates are the simultaneous eigenvalue tuples of these matrices.
\end{theorem}
\begin{proof}
The algebra $A_F$ is a nonzero finite-dimensional commutative $K$-algebra of dimension $\mu$. It has finitely many maximal ideals and a product decomposition into its Artinian localizations:
\[
 A_F\simeq\prod_{\mathfrak m\in\Max A_F}(A_F)_{\mathfrak m}.
\]
Every residue field is $K$. By \cref{thm:weaknull}, its maximal ideals correspond exactly to points of $Z_F$. The dimensions of the factors sum to $\mu$, and by \cref{prop:completion} these dimensions equal the formal local multiplicities \eqref{eq:localmultiplicity}. This proves the count.

In a local factor associated with $a$, multiplication by $z_j$ is $a_j$ times the identity plus a nilpotent operator. Simultaneous triangularization of the commuting multiplication operators gives precisely these tuples on the diagonal, with the corresponding local lengths. Conversely, a common eigenvector in the regular representation determines a character: if $z_jv=a_jv$, all polynomials act by evaluation, and the commuting finite algebra has the resulting maximal ideal. The evaluation description already identifies this character with a point of $Z_F$.
\end{proof}
The theorem protects total \emph{length}. It neither requires nor implies that the number of distinct points stays constant. A point of length $\mu$ can split into $\mu$ simple points or into intermediate nonreduced clusters.

\begin{proposition}[The simple-root criterion]\label{prop:jacobian}
The following conditions are equivalent:
\begin{enumerate}
\item Every point of $Z_F$ has local multiplicity one.
\item $J_F\sh(a)\ne0$ for all $a\in Z_F$.
\item $J_F$ is a unit in $A_F$.
\item The multiplication matrix $m_{J_F}$ has nonzero determinant.
\end{enumerate}
When they hold, $Z_F$ has exactly $\mu$ distinct points.
\end{proposition}
\begin{proof}
In a local factor at $a$, the completed ambient ring is $K[[h_1,\ldots,h_n]]$. The dimension of the cotangent space of its quotient by $F$ is $n$ minus the rank of the Jacobian matrix at $a$. An invertible Jacobian therefore makes this cotangent space zero, and Nakayama's lemma makes the local maximal ideal zero. Conversely, a local quotient of length one has zero cotangent space and full Jacobian rank. This proves the first equivalence. An element of an Artinian algebra is a unit exactly when it is absent from every maximal ideal, proving the second. Multiplication by an element of a finite-dimensional algebra is invertible exactly when that element is a unit, proving the last equivalence.
\end{proof}

\begin{theorem}[Workspace base change and full surcomplex exhaustiveness]\label{thm:basechange}
Suppose $\Gamma\subset\Gamma'$ are nonzero divisible ordered subgroups of $\No$, and the data of \cref{thm:conservation} lie in $\RR_n(\Gamma)$. Write $K'=\C((t^{\Gamma'}))$. Then the natural map is an isomorphism
\begin{equation}\label{eq:basechange}
 A_F(\Gamma)\otimes_KK'\ \xrightarrow{\ \sim\ }\ A_F(\Gamma').
\end{equation}
In the prescribed polynomial basis the multiplication matrices on both sides are identical. The zero set in $\mm_{K'}^n$ is exactly the original zero set in $\mm_K^n$, and its local multiplicities are unchanged.

Consequently, evaluating the same coherent system on the entire surcomplex monad introduces no additional roots. The finite count \eqref{eq:rootcount} is valid there as well.
\end{theorem}
\begin{proof}
The construction \eqref{eq:normalformoperator} depends only on the ordinary contraction and on the existing perturbation coefficients. Enlarging the exponent group does not change any coefficient calculation. By \cref{thm:conservation}, the same $b_1,\ldots,b_\mu$ are a basis on each side of \eqref{eq:basechange}; the natural map sends one basis to the other and is therefore an isomorphism of algebras.

The characteristic polynomial of each multiplication matrix has coefficients in $K$. All its roots already belong to $K$, since $K$ is algebraically closed. Any coordinate of a new common zero in $K'$ would have to be one of these roots. Its coordinates are therefore in $K$, and evaluation agrees with the original evaluation, so no new point appears. Base change of the Artinian local factors preserves their dimensions and hence their lengths.

Finally, the finitely many coordinates of any proposed surcomplex point have set-sized supports. Enlarge $\Gamma$ by those supports and take the divisible subgroup they generate. This gives a set-sized $\Gamma'$ to which the preceding argument applies. Thus every point in the full class monad was already in the original workspace.
\end{proof}
This is a stronger conclusion than placing each separately constructed root into some larger field. One fixed workspace containing the input already contains \emph{all} roots of the isolated cluster. Their supports need not lie in $T^*$: fractional scales such as $t^{1/2}$ may appear in individual roots even when matrix coefficients use only integral powers. Divisibility of $\Gamma$ accommodates these scales.

\section{A multivariate residue functional that survives collisions}\label{sec:residues}
\subsection{The classical input and the contour convention}
For an ordinary isolated complete intersection $f$ in a small neighborhood of zero, choose a standard Grothendieck residue cycle
\[
 \mathcal C_f=\{z:|f_1(z)|=\epsilon_1,\ldots,
                    |f_n(z)|=\epsilon_n\}
\]
in a suitable sufficiently small finite-map neighborhood. The radii are chosen so this is a compact regular real $n$-cycle, with orientation specified by
$d\arg f_1\wedge\cdots\wedge d\arg f_n>0$. Standard equivalent residue cycles give the same integral. Its normalization is
\begin{equation}\label{eq:classicalresidue}
 \lambda_f(g)=\frac1{(2\pi i)^n}
 \int_{\mathcal C_f}\frac{g\,dz_1\wedge\cdots\wedge dz_n}
                            {f_1\cdots f_n}.
\end{equation}

We use three classical residue facts: this functional annihilates $(f)$; its pairing on the finite complete-intersection algebra is nondegenerate; and it satisfies
\[
 \Tr(m_g)=\lambda_f(g\det Df).
\]
For a finite cluster of simple zeros the functional is the sum of $g(a)/\det Df(a)$. These are ordinary Grothendieck residue and complete-intersection duality statements, not new results of this article \cite[Chapter 5]{GH}\cite{SS}. The same statements hold for sufficiently small finite-dimensional holomorphic perturbation families, summing over the zeros in the fixed neighborhood; their total multiplicity is constant. We explain the comparison with such families in the proof of \cref{thm:traceresidue}.

For a uniformly supported Hahn family of differential forms on a compact ordinary cycle, define its $H$-integral coefficientwise. This is a coefficient functional; the cycle is not being treated as a fine-continuous parametrized object in $\SC^n$.

\subsection{Construction, descent, and duality}
On a neighborhood of $\mathcal C_f$ every $f_i$ is nonzero. Therefore
\begin{equation}\label{eq:residueexpansion}
 \frac1{F_i}=\frac1{f_i}\sum_{k\ge0}
          \left(-\frac{E_i}{f_i}\right)^k
\end{equation}
is a legitimate common-support Hahn expansion there. Define
\begin{equation}\label{eq:hahnresidue}
 \Lambda_F(G)=\frac1{(2\pi i)^n}
 \HInt_{\mathcal C_f}\frac{G\,dz_1\wedge\cdots\wedge dz_n}
                              {F_1\cdots F_n}.
\end{equation}
Here the differential form is expanded by \eqref{eq:residueexpansion} before integrating its ordinary coefficients.

\begin{theorem}[Support-controlled residue duality]\label{thm:residueduality}
For the data of \cref{thm:conservation}, \eqref{eq:hahnresidue} is well defined and is independent of the chosen sufficiently small ordinary residue cycle. It descends to a $K$-linear functional $\Lambda_F:A_F\to K$, and
\begin{equation}\label{eq:residuesupport}
 \supp\Lambda_F(G)\subset\supp G+T^*.
\end{equation}
The pairing
\begin{equation}\label{eq:pairing}
 A_F\times A_F\longrightarrow K,\qquad
 (a,b)\longmapsto\Lambda_F(ab)
\end{equation}
is perfect. In the prescribed polynomial basis its Gram matrix $W_F$ has support in $T^*$ and satisfies
\[
 W_F=W_0+\text{positive-support terms},\qquad
 W_0=(\lambda_f(b_ib_j))_{i,j},\qquad\det W_0\ne0.
\]
All these assertions remain valid when roots collide and the algebra becomes nonreduced.
\end{theorem}
\begin{proof}
Expanding all denominators gives coefficients supported in $\supp G+T^*$, with finitely many contributions at each exponent. On the cycle each contribution is an ordinary differential form, so the integral exists coefficientwise and has the stated support. Its individual terms are ordinary Grothendieck residue expressions with denominators $f_1^{k_1+1}\cdots f_n^{k_n+1}$. Their ordinary cycle independence proves independence of \eqref{eq:hahnresidue}, including passage to smaller common neighborhoods.

For $G=F_jH$, cancel the factor $F_j$ before expanding. Every resulting coefficient is an ordinary form with no $f_j$ in its denominator. Its integral is zero: after adjoining a factor $f_j$ to numerator and denominator it is a residue for the complete intersection of the indicated powers of the $f_i$, with numerator in its defining ideal. Thus $\Lambda_F$ annihilates $(F_1,\ldots,F_n)$. Scalar Hahn linearity follows from finite convolution at every exponent.

For the Gram matrix, each $b_ib_j$ is ordinary. At exponent zero the errors in the denominator expansion disappear, leaving exactly the classical matrix $W_0$. The ordinary residue pairing is nondegenerate. Consequently the determinant of $W_F$ has nonzero coefficient $\det W_0$ at exponent zero, so $W_F$ is invertible over $K$. Since the classes of the $b_i$ form a basis of $A_F$, the pairing is perfect. The argument does not use nonvanishing of a Jacobian at any root, which proves the collision assertion.
\end{proof}

\begin{corollary}[A support-controlled constant residue frame]\label{cor:residueframe}
Let $C=W_0^{-1}(W_F-W_0)$. The matrix series
\begin{equation}\label{eq:residueframe}
 S=(I+C)^{-1/2}
   =\sum_{k\ge0}\binom{-1/2}{k}C^k
\end{equation}
is strongly summable, has support in $T^*$, and satisfies
\[
 S^{\mathsf T}W_FS=W_0.
\]
Thus the residue bilinear space has a basis congruent to the prescribed ordinary basis in which its matrix is exactly the ordinary residue matrix, even on the collision locus.
\end{corollary}
\begin{proof}
The matrix $C$ has positive support. The support lemma makes every coefficient of \eqref{eq:residueframe} a finite sum; ordinary one-variable formal power-series identities in the single matrix $C$ give $S^2(I+C)=I$. Symmetry of $W_F$ and $W_0$ gives $C^{\mathsf T}W_0=W_0C$, hence $S^{\mathsf T}W_0=W_0S$. Since $W_F=W_0(I+C)$ and $S$ commutes with $C$, the displayed isometry follows.
\end{proof}
This identifies bilinear spaces, not algebras. It does not erase the nonreduced algebraic structure of a collision. In particular, a perfect residue pairing should not be confused with the usually degenerate trace pairing of a nonreduced finite algebra.

\subsection{Trace and the actual moving roots}
\begin{theorem}[Trace--Jacobian identity and moving-root formula]\label{thm:traceresidue}
For every $G\in A_F$,
\begin{equation}\label{eq:traceidentity}
 \Tr_{A_F/K}(m_G)=\Lambda_F(GJ_F).
\end{equation}
In particular,
\begin{equation}\label{eq:residuecount}
 \Lambda_F(J_F)=\mu.
\end{equation}
If all points of $Z_F$ are simple, then
\begin{equation}\label{eq:movingresidues}
 \Lambda_F(G)=\sum_{a\in Z_F}
       \frac{G\sh(a)}{J_F\sh(a)}.
\end{equation}
The support bound \eqref{eq:residuesupport} applies to the total in \eqref{eq:movingresidues}, even when its individual summands have negative valuation.
\end{theorem}
\begin{proof}
We first justify \eqref{eq:traceidentity} coefficientwise, taking care not to assume ordinary convergence of the Hahn perturbation. At any fixed output exponent, the multiplication matrices in \cref{thm:conservation} and the contour expansion \eqref{eq:hahnresidue} use only finitely many words in perturbation coefficients. This follows from the support bounds and \cref{lem:neumann}. Label all the finitely many coefficient functions occurring in these words, and introduce independent ordinary complex parameters for them. We obtain a finite-dimensional holomorphic perturbation
\[
 f(z)+\sum_{j=1}^{q}u_jg_j(z),\qquad u\in\C^q.
\]
For small $u$, ordinary analytic conservation of number gives a finite cluster of total multiplicity $\mu$ in the chosen neighborhood. Its quotient algebras form a locally free holomorphic family of rank $\mu$, and the selected $b_j$ are a basis after shrinking the parameter neighborhood. One standard justification is that the relative zero space is finite over the parameter domain, is a complete intersection and hence Cohen--Macaulay, and is therefore flat over this smooth parameter base; the basis assertion follows from the ordinary fiber at $u=0$. These are classical finite analytic map and complete-intersection facts \cite{GR,GH,Stacks}.

In this finite-parameter family, multiplication matrices and the fixed-cycle residues are holomorphic functions of $u$. The classical trace--Jacobian identity holds for each sufficiently small $u$. Its Taylor series identity therefore holds formally. The Taylor series of the ordinary normal form is exactly \eqref{eq:normalformoperator}, now interpreted as a formal series in the finitely many $u_j$: both kill the deformed Koszul image and are the identity on the prescribed basis. Uniqueness follows also directly from \cref{thm:conjugacy}, using ordinary parameter degree to form its formal inverse.

Replace the parameter attached to a coefficient of exponent $\gamma$ by $t^\gamma$. For the target exponent under consideration there are only finitely many contributing words, so this substitution proves that coefficient of \eqref{eq:traceidentity}. This argument works at every exponent, including those in non-Archimedean value groups. It first proves the identity for ordinary representatives and then for arbitrary Hahn $G$ by the same finite-contribution argument, or simply by $K$-linearity on the finite basis.

Taking $G=1$ gives \eqref{eq:residuecount}. If every point is simple, \cref{prop:jacobian} makes $J_F$ invertible in $A_F$. Apply \eqref{eq:traceidentity} to $G/J_F$. The reduced finite algebra is $K^{Z_F}$, and multiplication acts diagonally with eigenvalues $G\sh(a)/J_F\sh(a)$. Its trace is their sum, proving \eqref{eq:movingresidues}. The support bound follows from \cref{thm:residueduality}.
\end{proof}

\begin{remark}[What is and is not being localized]
At a multiple point, \eqref{eq:movingresidues} is not a valid formula with a zero denominator. The always-valid statement is the functional on its finite local algebra. The Artinian product decomposition restricts $\Lambda_F$ to a perfect functional on each local factor, and their sum is the total functional. Its trace--Jacobian identity still holds. This gives an algebraic local residue interpretation without inventing small fine-topological integration cycles around individual nonstandard points.
\end{remark}

\section{A length-five collision with three scale regimes}\label{sec:example}
This example makes the finite algebra, actual roots, collision equation, and residue pairing explicit. Let $s,u\in\mm_K$ and consider
\begin{equation}\label{eq:exampleF}
 F_1=x^2-y^3-s,\qquad F_2=xy-u.
\end{equation}
At $s=u=0$, the ordinary zero is isolated. Its local algebra has basis
\begin{equation}\label{eq:examplebasis}
 \bvec=(1,y,y^2,y^3,x),\qquad \mu=5.
\end{equation}
For example, the leading monomials of a local polynomial reduction are $x^2,xy,y^4$, leaving exactly these five monomials. Alternatively the relations below directly span the quotient and their commuting matrices prove independence.

\subsection{Multiplication matrices and elimination}
The defining relations imply
\begin{equation}\label{eq:examplerelations}
 xy=u,\qquad x^2=s+y^3,\qquad y^4=ux-sy.
\end{equation}
With columns representing images of the ordered basis \eqref{eq:examplebasis}, multiplication is given by
\begin{equation}\label{eq:examplematrices}
 M_x=\begin{pmatrix}
 0&u&0&0&s\\
 0&0&u&0&0\\
 0&0&0&u&0\\
 0&0&0&0&1\\
 1&0&0&0&0
 \end{pmatrix},\qquad
 M_y=\begin{pmatrix}
 0&0&0&0&u\\
 1&0&0&-s&0\\
 0&1&0&0&0\\
 0&0&1&0&0\\
 0&0&0&u&0
 \end{pmatrix}.
\end{equation}
They satisfy the exact polynomial identities
\begin{equation}\label{eq:matrixidentities}
 M_xM_y=M_yM_x=uI,
 \qquad M_x^2-M_y^3=sI.
\end{equation}
Starting with the vector for $1$, the matrices generate the five basis vectors, so they also give a faithful five-dimensional realization of the algebra described by \eqref{eq:examplerelations}. The conservation theorem supplies the same conclusion for the analytic monad quotient.

Their characteristic polynomials are
\begin{equation}\label{eq:characteristicexample}
 \chi_y(Y)=Y^5+sY^2-u^2,\qquad
 \chi_x(X)=X^5-sX^3-u^3.
\end{equation}
When $u\ne0$, every root of $\chi_y$ is nonzero and determines one solution by $x=u/y$. In all cases the total local length is five. The polynomial equations are elimination tools, but a chosen coordinate need not distinguish all points.

\subsection{The actual collision equation}
The Jacobian is
\begin{equation}\label{eq:exampleJ}
 J=\det\begin{pmatrix}2x&-3y^2\\y&x\end{pmatrix}
   =2x^2+3y^3=2s+5y^3\quad\text{in }A_F.
\end{equation}
Direct calculation gives
\begin{equation}\label{eq:collision}
 \det(m_J)=3125u^6-108s^5.
\end{equation}
Hence the system has five simple points exactly when the right side of \eqref{eq:collision} is nonzero. This is the collision equation of the \emph{whole algebra}, not merely a discriminant of a projected coordinate.

Indeed,
\begin{equation}\label{eq:projecteddiscriminant}
 \disc_Y(Y^5+sY^2-u^2)=u^2(3125u^6-108s^5).
\end{equation}
The extra factor $u^2$ is a projection artifact. At $u=0$, $s\ne0$, the system has two points $(\pm\sqrt{s},0)$ and three points $(0,y)$ with $y^3=-s$. All five are simple, even though the polynomial in $y$ has a double root at zero. This illustrates why simultaneous multiplication matrices and the full Jacobian provide better finite-map data than a single eliminant.

A nontrivial collision is explicit. Choose any $a\in\mm_K^\times$, choose $\eta\in\C$ with $\eta^2=-3/2$, and take
\[
 s=-\frac52a^6,\qquad u=\eta a^5.
\]
Then $(x,y)=(\eta a^3,a^2)$ is a zero with $J=0$, and \eqref{eq:collision} vanishes. In the $y$-eliminant this is exactly a double root: its second derivative there is $15a^6\ne0$. After writing $y=a^2Y$, the eliminant is proportional to $(Y-1)^2(2Y^3+4Y^2+6Y+3)$; its greatest common divisor with its derivative is $Y-1$. Thus this specialization has one point of length two and three other simple points. At $s=u=0$ all five units of length are concentrated at the origin.

\subsection{Three infinitesimal scale regimes}
Take positive exponents $\alpha,\beta\in\Gamma$ and set
\[
 s=t^\alpha,\qquad u=t^\beta.
\]
The relative sizes are governed by $5\alpha$ and $6\beta$. All roots are simple: in the equal case the two nonzero leading coefficients in \eqref{eq:collision} subtract to $3125-108=3017$, not zero.

\begin{center}
\begin{tabular}{@{}L{0.23\textwidth}L{0.12\textwidth}L{0.23\textwidth}L{0.23\textwidth}@{}}
\toprule
Regime & Number & $v(y)$ & $v(x)$\\
\midrule
$5\alpha<6\beta$ & $2$ & $\beta-\alpha/2$ & $\alpha/2$\\[0.3em]
 & $3$ & $\alpha/3$ & $\beta-\alpha/3$\\[0.3em]
$5\alpha>6\beta$ & $5$ & $2\beta/5$ & $3\beta/5$\\[0.3em]
$5\alpha=6\beta$ & $5$ & $2\beta/5$ & $3\beta/5$\\
\bottomrule
\end{tabular}
\end{center}
We prove these assertions directly by rescaling, rather than merely displaying a Newton polygon.

Suppose first that $5\alpha<6\beta$. Substitute $y=t^{\beta-\alpha/2}Y$ in $\chi_y=0$ and divide by $t^{2\beta}$. The equation becomes
\[
 Y^2-1+t^{3\beta-5\alpha/2}Y^5=0.
\]
The error exponent is positive. Each of the two simple ordinary roots $Y=\pm1$ lifts uniquely to a root in its monad, by degree-one preparation or the simple-root case of \cref{thm:rootcount}. This gives the first group. Next put $y=t^{\alpha/3}Y$ and divide by $t^{5\alpha/3}$, obtaining
\[
 Y^5+Y^2-t^{2\beta-5\alpha/3}=0.
\]
Its three nonzero simple leading roots satisfy $Y^3=-1$. They lift to the second group. The valuations in the two groups differ, and together they account for all five roots. Since $x=t^\beta/y$, the stated $x$-valuations follow.

For $5\alpha>6\beta$, put $y=t^{2\beta/5}Y$ and divide by $t^{2\beta}$. This gives
\[
 Y^5-1+t^{\alpha-6\beta/5}Y^2=0.
\]
Its five ordinary fifth roots of unity are simple and lift uniquely. In the equal case the same substitution gives the ordinary polynomial $Y^5+Y^2-1$. It has no zero root and is squarefree, since its discriminant is $3017$. Its five ordinary roots give the five solutions at that scale. This proves the table in every divisible ordered value group, not just for real rational exponents.

\subsection{Residues without singular denominators}
In the basis \eqref{eq:examplebasis}, the total residue functional is
\begin{equation}\label{eq:examplelambda}
 \Lambda(1)=\Lambda(y)=\Lambda(y^2)=\Lambda(x)=0,
 \qquad \Lambda(y^3)=1.
\end{equation}
One way to identify it first on the simple locus with $u\ne0$ is to use the $y$-eliminant $q(Y)=Y^5+sY^2-u^2$. At a root,
\[
 q'(y)=yJ.
\]
The classical coefficient identity for a monic degree-five polynomial gives
\[
 \sum_{q(y)=0}\frac{y^k}{q'(y)}=0\ (0\le k<4),
 \qquad \sum_{q(y)=0}\frac{y^4}{q'(y)}=1.
\]
These formulas and $x=u/y$ give \eqref{eq:examplelambda}. The fixed-cycle residue depends holomorphically on the ordinary parameters near $(0,0)$; equality on the open simple locus therefore proves the identities throughout the parameter germ, and hence at every infinitesimal Hahn specialization. They may also be obtained from the ordinary residue expansion directly.

The Gram matrix is
\begin{equation}\label{eq:exampleGram}
 W=\begin{pmatrix}
 0&0&0&1&0\\
 0&0&1&0&0\\
 0&1&0&0&0\\
 1&0&0&-s&0\\
 0&0&0&0&1
 \end{pmatrix},\qquad \det W=1.
\end{equation}
Unlike \eqref{eq:collision}, this determinant never vanishes. The residue pairing remains perfect at every collision, including the length-five point at the origin. By contrast its trace Gram matrix is $Wm_J$, so
\[
 \det\bigl(\Tr(m_{b_ib_j})\bigr)_{i,j}=3125u^6-108s^5,
\]
and does degenerate on the collision locus. Also
\[
 \Lambda(J)=\Lambda(2s+5y^3)=5,
\]
which is the exact count. On the simple locus this gives, among others,
\begin{equation}\label{eq:examplecancellation}
 \sum_{F(a)=0}\frac1{J(a)}=0,
 \qquad
 \sum_{F(a)=0}\frac{y(a)^3}{J(a)}=1,
 \qquad
 \sum_{F(a)=0}1=5.
\end{equation}
For monomial parameters in the scale table, each $J(a)$ is a nonzero infinitesimal, so every summand in the first sum has infinite modulus. Their cancellation is nevertheless exact. The total residue is computed by a fixed nonsingular matrix rather than by subtracting poorly separated pole contributions.

\section{Beyond polynomial and grid-supported perturbations}\label{sec:nonpolynomial}
The finite example is not the limit of the theorem's scope. Consider the common-domain entire coefficient families
\begin{align}\label{eq:nonpolynomial}
 F_1(x,y)&=x^2-y^3-
      \sum_{m\ge1}t^{\,2-1/m}\exp(m!x),\\
 F_2(x,y)&=xy-t^{\sqrt2}-t^3\sin(x+y).\nonumber
\end{align}
Choose a divisible workspace containing the real exponents involved. The support
\[
 T=\{2-1/m:m\ge1\}\cup\{\sqrt2,3\}
\]
is positive and well ordered. Its subset $\{2-1/m\}$ has a finite upper cut at $2$, while its denominators are unbounded; the construction is not restricted to a discrete one-generator grid. All coefficient functions are entire, so a common ordinary domain is immediate despite their very large growth.

\begin{corollary}[A nonpolynomial length-five cluster]\label{cor:nonpolynomial}
The system \eqref{eq:nonpolynomial} has finitely many common zeros in the full surcomplex monad of $(0,0)$, of total local multiplicity five. Its quotient algebra has the exact basis $(1,y,y^2,y^3,x)$. Every entry of its two multiplication matrices and of its residue Gram matrix has support contained in $T^*$. The residue Gram matrix is invertible, and the total residue of its Jacobian equals five.
\end{corollary}
\begin{proof}
Its ordinary leading system is exactly $(x^2-y^3,xy)$, of multiplicity five. Apply \cref{thm:conservation,thm:rootcount,thm:basechange,thm:residueduality,thm:traceresidue}. Each input to a multiplication or Gram matrix entry is ordinary, giving the support bound $T^*$.
\end{proof}
This corollary does not assert that the explicit polynomial matrices \eqref{eq:examplematrices} remain valid for the nonpolynomial system; they do not. It asserts that the same \emph{basis size and basis representatives} survive, and that the general normal-form recursion computes the new entries. Nor does it assert that all five roots are simple without computing the new Jacobian determinant in that algebra.

The proof is unaffected by replacing $T$ with a positive well-ordered set in a higher-rank subgroup of $\No$, or by using several distinct Archimedean classes of scales. The only summation requirements are well ordering and finite contribution at each exponent. No ordinary bound on all coefficient derivatives is required.

\section{The scope of the resolution}\label{sec:scope}
\subsection{What has been established}
For the several-variable direction raised in \cite{Source}, the common-domain category now has parameterized preparation and division, a Noetherian point-detecting germ ring, finite projections, exact isolated-cluster multiplicity, support-controlled normal forms and matrices, workspace independence, and a multivariate moving-root residue formula. The central invariant is the finite quotient algebra. Individual roots and ordinary-contour coefficients are two ways to read it, not competing definitions of multiplicity.

The proof separates three resources. Ordinary complex analysis supplies one fixed-domain Koszul contraction and the classical residue duality input. Hahn support calculus transports those constructions without topological iteration limits. Algebraic closedness and the monad Nullstellensatz identify the resulting finite algebra with every actual surcomplex point of the cluster. Omitting any one of these steps would leave a gap between formal manipulations and the stated zero-count theorem.

\subsection{Hypotheses that cannot be dropped silently}
\paragraph{An isolated complete intersection.}
The leading system must have finite local colength. A leading system such as $(xy,xy)$ in two variables has a positive-dimensional zero set, so no finite number $\mu$ is available. Its perturbations may create very different geometries. The theorem does not claim conservation of a nonexistent finite length.

\paragraph{Positive perturbation support.}
Replacing $x$ by $x-1$ changes the ordinary leading equation and removes its zero from the origin's monad. An error with an exponent below zero can dominate the leading system entirely. Positive support is the mechanism making $1+\delta h$ invertible with the stated support certificate.

\paragraph{One common ordinary neighborhood.}
Individually holomorphic coefficient germs with radii tending to zero need not be coherent data. The fixed-domain construction is therefore not an optional analytic convenience. Conversely, coefficients may have arbitrarily large norms on that common domain; no unmentioned boundedness hypothesis is used.

\paragraph{A set-sized workspace before passing to the class.}
Noetherianity and finite-dimensional linear algebra are asserted for the set-sized rings $\RR_n(\Gamma)$. The passage to the full surreal class is made by the explicit base-change argument for each finite cluster, not by treating a proper class as an ordinary Noetherian ring.

\paragraph{The chosen coherent analytic category.}
These results do not classify arbitrary fine-locally analytic functions on the class plane. In particular, they do not settle the supplied manuscript's separate questions about essential singularities, resurgent sectorial summation, or transcendental all-scale atlases. Those involve compatibility of different expansion structures rather than just a fixed positive-support deformation of an ordinary finite germ.

\subsection{Further precise targets}
The next extension is a relative version over a positive-dimensional ordinary parameter space: construct normal forms as modules over the \emph{coherent parameter algebra}, rather than merely as complex-linear splittings on the total space. The proof above deliberately does not claim this stronger relative linearity. A useful formulation would start with a finite flat ordinary family and a chosen locally free quotient bundle, and ask for support-controlled lifting of that bundle and its residue duality.

Another target is to replace an isolated complete intersection by a finite Cohen--Macaulay module with a selected free resolution. The positive-conjugacy argument in \cref{thm:conjugacy} already applies to a complex with homology only in degree zero, provided a perturbation of its differential is supplied and still squares to zero. The genuinely new problem is then to construct the required compatible perturbation of all syzygies, rather than just of the first equations. For complete intersections, the Koszul construction supplies it automatically.

These targets distinguish what is proved here from what a broader several-variable surcomplex analytic geometry would still require.

\appendix
\section{A proof and support audit}\label{app:audit}
The dependency chain is
\begin{gather*}
 \text{Hahn--Neumann lemma}
 \ \Longrightarrow\ \text{positive operator inversion},\\
 \text{positive operator inversion}
 \ \Longrightarrow\ \text{Koszul conjugacy}
 \ \Longrightarrow\ \text{fixed-basis conservation}.
\end{gather*}
Preparation supplies the Noetherian and point-identification route. Classical residue theory supplies the leading perfect form and the finite-parameter trace identity; support calculus then gives the Hahn extension.

\begin{center}\small
\begin{longtable}{@{}L{0.39\textwidth}L{0.52\textwidth}@{}}
\toprule
Construction & Controlling support or finite-dimensional fact\\
\midrule
\endhead
Evaluation at $c+\eps$ & $\supp F+(\bigcup_i\supp\eps_i)^*$\\[0.4em]
Infinitesimal translation by $a$ & $\supp F+(\bigcup_i\supp a_i)^*$\\[0.4em]
Ordinary linear operator applied coefficientwise & A subset of the input support\\[0.4em]
Preparation factors & $T^*$ for the normalized error support $T$\\[0.4em]
Division by $P_0+A$ & $\supp G+(\supp A)^*$\\[0.4em]
Koszul inverse $(1+\delta h)^{-1}$ & $\supp G+T^*$\\[0.4em]
Normal forms and division certificates & $\supp G+T^*$, on one ordinary domain\\[0.4em]
Multiplication matrices in an ordinary basis & $T^*$; reduction gives ordinary nilpotent matrices\\[0.4em]
Total residue of $G$ & $\supp G+T^*$\\[0.4em]
Residue Gram matrix and constant frame & $T^*$; the leading Gram determinant is nonzero\\[0.4em]
Individual root coordinates & Algebraic over $K$; not necessarily supported in $T^*$\\[0.4em]
Workspace enlargement & Same finite basis, same matrices, algebraic closedness of the original $K$\\
\bottomrule
\end{longtable}
\end{center}

\subsection{Finite coefficients do not mean finite initial segments}
A well-ordered positive support may have infinitely many exponents below a specified bound. Our proofs never require those initial segments to be finite. They require finitely many words contributing to \emph{one exact exponent}. For example, if $T=\{2-1/m:m\ge1\}$, there are infinitely many members below $2$, but the coefficient of a given Hahn product is still a finite convolution. This distinction is also what makes \cref{ex:ranktwo} legitimate.

\subsection{Where classical theorems enter}
The article gives proofs of the Hahn extensions and their support claims, not new proofs of the entire ordinary Oka--Cartan theory. Its classical inputs are ordinary Weierstrass preparation, coherence and Cartan B on Stein domains, the regular-sequence property of a system of parameters in a regular local ring, standard integral-extension and Artinian decomposition theorems, and ordinary complete-intersection residue duality. Each is explicitly used at its corresponding proof step and cited there.

The finite-parameter comparison in \cref{thm:traceresidue} does not approximate a Hahn family by a convergent sequence of finite truncations. Rather, it verifies one coefficient using a finite number of independent ordinary parameters. The support lemma then permits those coefficient identities to be assembled. No interchange of a fine-topological limit with an integral is involved.

\section{Exact computational checks}\label{app:checks}
The accompanying \texttt{verify\_examples.py} uses exact symbolic arithmetic to check
\begin{enumerate}
\item both defining matrix relations and the two characteristic polynomials in \cref{sec:example};
\item the full Jacobian determinant, the projected discriminant, and their distinct collision information;
\item the residue Gram matrix, its determinant, the trace--Jacobian identity on a basis, and the matrix realization of the residue functional;
\item the explicit nontrivial double-root specialization and its second derivative;
\item the exponents and signs in the three scale rescalings, and the discriminant $3017$ of the balanced residual quintic.
\end{enumerate}
The report is generated by running that script, not by manually asserting that it passed. These finite exact checks verify the displayed example and detect sign or indexing mistakes. They are not a formal proof-assistant verification of the general theorems. The latter are supported by the arguments and support audits in the article.

\section{Notation}
\begin{center}\small
\begin{longtable}{@{}L{0.26\textwidth}L{0.65\textwidth}@{}}
\toprule
Symbol & Meaning\\
\midrule
\endhead
$\SC=\No[i]$ & Surcomplex class field\\[0.3em]
$t^\gamma=\omega^{-\gamma}$ & Normal-form monomial; positive exponent means infinitesimal\\[0.3em]
$\Gamma$, $K_\Gamma$ & Nonzero divisible set-sized value group and its complex Hahn field\\[0.3em]
$\mm_K$ & Positive-valuation elements of $K$, together with zero\\[0.3em]
$\HH_\Gamma(U)$ & Hahn families with all ordinary holomorphic coefficients on $U$\\[0.3em]
$\RR_n(\Gamma)$ & Germ ring of common-domain Hahn families at the ordinary origin\\[0.3em]
$T^*$ & Positive-support additive monoid, including the empty sum\\[0.3em]
$A_0$, $A_F$ & Ordinary local algebra and its Hahn-perturbed monad algebra\\[0.3em]
$\mu$, $\mu_a(F)$ & Total ordinary multiplicity and local deformed multiplicity\\[0.3em]
$d_f$, $\delta$, $h$ & Ordinary Koszul differential, its positive perturbation, and a fixed contraction\\[0.3em]
$N(G)$, $Q(G)$ & Normal-form vector and division certificate\\[0.3em]
$M_j$ & Multiplication matrix of the coordinate $z_j$\\[0.3em]
$J_F$ & Determinant of the coefficientwise Jacobian matrix of $F$\\[0.3em]
$\Lambda_F$, $W_F$ & Total Hahn residue functional and its Gram matrix\\
\bottomrule
\end{longtable}
\end{center}

\begin{thebibliography}{99}
\addcontentsline{toc}{section}{References}

\bibitem{Source}
\emph{Surcomplex Analysis: Infinitesimal Calculus, Hahn-Coherent Holomorphy, and Contour Theory}.
User-supplied unpublished manuscript, supplied file
\texttt{surcomplex\_analysis(2).tex}, dated September 21, 2026.
The present article uses its Hahn-coherent framework and addresses its concluding several-variable research direction. No authorship or published priority is inferred from the file.

\bibitem{CLR}
Raf Cluckers, Leonard Lipshitz, and Zachary Robinson,
``Real closed fields with nonstandard and standard analytic structure,''
\emph{Journal of the London Mathematical Society} (2) \textbf{78} (2008), 198--212.
\href{https://doi.org/10.1112/jlms/jdn024}{doi:10.1112/jlms/jdn024};
\href{https://arxiv.org/abs/math/0610666}{arXiv:math/0610666}.

\bibitem{CL}
Raf Cluckers and Leonard Lipshitz,
``Fields with analytic structure,''
\emph{Journal of the European Mathematical Society} \textbf{13} (2011), 1147--1223.
\href{https://doi.org/10.4171/JEMS/278}{doi:10.4171/JEMS/278};
\href{https://arxiv.org/abs/0908.2376}{arXiv:0908.2376}.
The arbitrary-support, common-radius examples are in \S3.1 and Example 3.3(6) of the arXiv text; related normalization and Nullstellensatz results appear in \S5.2.

\bibitem{Crainic}
Marius Crainic,
``On the perturbation lemma, and deformations,'' 2004.
\href{https://arxiv.org/abs/math/0403266}{arXiv:math/0403266}.
Our contraction convention is $dh+hd=1-ip$; the formulas in \cref{thm:conjugacy} are proved with that sign convention.

\bibitem{GR}
Robert C. Gunning and Hugo Rossi,
\emph{Analytic Functions of Several Complex Variables},
Prentice-Hall, 1965; reprinted by AMS Chelsea Publishing, 2009.
Ordinary Weierstrass theory, coherent analytic sheaves, and Stein-space theorems.

\bibitem{GH}
Phillip Griffiths and Joseph Harris,
\emph{Principles of Algebraic Geometry}, Wiley, 1978; Wiley Classics Library edition, 1994.
Chapter 5, ``Residues.''
\href{https://doi.org/10.1002/9781118032527}{doi:10.1002/9781118032527}.

\bibitem{SS}
G\"unter Scheja and Uwe Storch,
``\"Uber Spurfunktionen bei vollst\"andigen Durchschnitten,''
\emph{Journal f\"ur die reine und angewandte Mathematik} \textbf{278/279} (1975), 174--190.
\href{https://doi.org/10.1515/crll.1975.278-279.174}{doi:10.1515/crll.1975.278-279.174}.

\bibitem{Neumann}
B. H. Neumann,
``On ordered division rings,''
\emph{Transactions of the American Mathematical Society} \textbf{66} (1949), 202--252.
Support and well-ordering results for generalized series.

\bibitem{Poonen}
Bjorn Poonen,
``Maximally complete fields,''
\emph{L'Enseignement Math\'ematique} (2) \textbf{39} (1993), 87--106.
\href{https://math.mit.edu/~poonen/papers/amsval.pdf}{Author's full text}.
In particular, the algebraic-closure result for divisible value group and algebraically closed residue field.

\bibitem{Stacks}
The Stacks Project Authors,
\emph{The Stacks Project}, chapter ``Commutative Algebra,''
sections on finite and integral extensions, regular local rings, Cohen--Macaulay rings, and flatness; consulted September 21, 2026.
\href{https://stacks.math.columbia.edu/tag/00GQ}{Tag 00GQ} gives the lying-over statement used here.

\end{thebibliography}
\end{document}
